\chapter{Motivação e Objetivos}\label{chap:problema}

Este capítulo apresenta a motivação que originou a proposta deste trabalho,
os objetivos geral e específicos a serem alcançados e as hipóteses de pesquisa formuladas.

\section{Motivação}

A interpretação de análises químicas e físicas de solo é uma etapa essencial
no planejamento da fertilidade e do manejo das culturas agrícolas. Trata-se de
um processo que exige elevado conhecimento técnico e atenção rigorosa, pois erros
de interpretação podem acarretar prejuízos financeiros significativos aos produtores,
seja pelo uso desnecessário de insumos, seja pela redução de produtividade decorrente
de correções insuficientes.

A complexidade desse processo advém da multiplicidade de variáveis envolvidas
--- pH, teores de macronutrientes e micronutrientes, capacidade de troca catiônica
(CTC), saturação por bases, textura do solo, entre outros \citep{Embrapa2011,
ManualCalagem2016} --- e da necessidade de correlacioná-las com as exigências
específicas de cada cultura e com as características edafoclimáticas de cada região.
Ferramentas computacionais que realizem essa correlação de forma estruturada ainda
são escassas no contexto do produtor e do técnico gaúcho, especialmente no que
diz respeito à integração entre recomendações de manejo e estimativas de aptidão
agrícola para culturas específicas \citep{Hofig2015, EmbrapaAptidao2025}.

Nesse contexto, o autor deste trabalho acumulou experiência profissional tanto
na prestação de assistência técnica a produtores rurais, com contato direto com
laudos de análise de solo e recomendações de correção e adubação, quanto no
atendimento bancário a carteiras de produtores rurais, em processos vinculados
à liberação de crédito rural. Nessa última atividade, observa-se que projetos
agrícolas submetidos às instituições financeiras podem apresentar inconsistências
nas recomendações técnicas, evidenciando a necessidade de instrumentos que apoiem
a validação dessas informações de maneira objetiva e rastreável.

Diante desse cenário, identificou-se a oportunidade de desenvolver o
\textbf{SIRAS} --- \textbf{S}istema de \textbf{I}nterpretação e \textbf{R}ecomendação
\textbf{A}gronômica de \textbf{S}olo ---, uma ferramenta computacional capaz de:
(i)~interpretar automaticamente dados de análises de solo; (ii)~gerar recomendações
técnicas de correção e manejo; e (iii)~estimar a aptidão agrícola de culturas
específicas em uma determinada área. O SIRAS poderia auxiliar técnicos e produtores
no processo de tomada de decisão, reduzir erros de interpretação e servir de apoio
na análise e validação de projetos agrícolas.

Iniciativas como o FertFacil, sistema desenvolvido e disponibilizado para
técnicos e produtores do RS, SC, MS e Paraguai, demonstram que a automação do
processo de recomendação de adubação e calagem é tecnicamente viável e tem gerado
resultados concretos no campo \citep{Tiecher2016}. No entanto, ferramentas dessa
natureza concentram-se na geração de recomendações de correção e adubação, não
contemplando a estimativa da aptidão agrícola da cultura escolhida frente às
condições edáficas da área. Em escala mais ampla, a Embrapa desenvolveu o Sistema
de Avaliação de Aptidão Agrícola das Terras \citep{EmbrapaAptidao2025}, voltado
a levantamentos exploratórios regionais, o que o torna inadequado para a avaliação
individualizada de uma área específica a partir dos dados de uma análise de solo.
Conforme apontado por \cite{Hofig2015}, esse sistema é mais adequado para
análises em âmbito regional e não indica práticas de manejo para os diferentes
tipos de utilização, sendo necessárias adaptações para estudos mais detalhados.
Constata-se, portanto, uma lacuna entre as ferramentas disponíveis: existem
sistemas que recomendam correção do solo, e existem métodos que avaliam aptidão
em escala regional, mas não há, no contexto das ferramentas acessíveis ao produtor
e ao técnico do RS, uma solução que integre ambas as perspectivas a partir de uma
análise individual de solo. Esse é o diferencial do SIRAS.

%-------------------------------------------------------------------
\section{Objetivo Geral}\label{sec:obj_geral}

Desenvolver e avaliar o SIRAS --- Sistema de Interpretação e Recomendação
Agronômica de Solo ---, um sistema especialista baseado em regras projetado
para interpretar dados de análises químicas e físicas de solo, gerar
recomendações técnicas de correção e manejo, e estimar a aptidão agrícola
para um conjunto abrangente de culturas de interesse regional, considerando
os critérios técnicos vigentes para o estado do Rio Grande do Sul.

%-------------------------------------------------------------------
\section{Objetivos Específicos}\label{sec:obj_espec}

\begin{itemize}[leftmargin=5em]

  \item \textbf{Revisar a literatura} sobre os principais parâmetros utilizados
    na interpretação de análises de solo, abrangendo aspectos químicos e físicos,
    suas correlações com o desempenho das culturas agrícolas contempladas pelo
    SIRAS, e os critérios de calagem e adubação vigentes para o Rio Grande do Sul;

  \item \textbf{Definir critérios técnicos} para as recomendações de correção
    e manejo do solo e para a estimativa de aptidão agrícola das culturas
    contempladas, com base nos manuais e normas regionais vigentes;

  \item \textbf{Projetar e implementar} o SIRAS, incluindo os módulos de
    entrada de dados de análise de solo, processamento das recomendações
    de correção, e estimativa de aptidão agrícola por cultura;

  \item \textbf{Validar o SIRAS} quanto à correção de suas saídas, medindo
    a taxa de concordância das recomendações geradas com valores de
    referência calculados de forma independente a partir dos critérios
    técnicos vigentes;

  \item \textbf{Avaliar a contribuição e a usabilidade do SIRAS}, comparando-o
    qualitativamente às ferramentas existentes para evidenciar seu diferencial
    e aferindo a experiência de uso junto ao público-alvo.

\end{itemize}

%-------------------------------------------------------------------
\section{Hipóteses}\label{sec:hipoteses}

Com base no problema identificado e nos objetivos propostos,
formulam-se as seguintes hipóteses de pesquisa:

\begin{itemize}

  \item \textbf{H\textsubscript{0.1}:} O SIRAS \textbf{não} é capaz de gerar
    recomendações de calagem e adubação com taxa de concordância igual ou
    superior a 80\% em relação aos valores de referência derivados dos
    critérios estabelecidos no Manual de Calagem e Adubação para os estados
    do RS e SC.

  \item \textbf{H\textsubscript{1.1}:} O SIRAS \textbf{é} capaz de gerar
    recomendações de calagem e adubação com taxa de concordância igual ou
    superior a 80\% em relação aos valores de referência derivados dos
    critérios estabelecidos no Manual de Calagem e Adubação para os estados
    do RS e SC.

  \item \textbf{H\textsubscript{0.2}:} O SIRAS \textbf{não} é capaz de estimar
    a aptidão edáfica das culturas contempladas com taxa de concordância igual
    ou superior a 80\% em relação à classificação de referência construída a
    partir das exigências agronômicas descritas na literatura técnica.

  \item \textbf{H\textsubscript{1.2}:} O SIRAS \textbf{é} capaz de estimar
    a aptidão edáfica das culturas contempladas com taxa de concordância igual
    ou superior a 80\% em relação à classificação de referência construída a
    partir das exigências agronômicas descritas na literatura técnica,
    diferenciando-se dos sistemas de avaliação regional existentes por operar
    a partir de uma análise individual de solo.

\end{itemize}

A confirmação de H\textsubscript{1.1} indicará que é viável automatizar,
ao menos parcialmente, o processo de interpretação de análises de solo e
geração de recomendações para o contexto do RS. A confirmação de
H\textsubscript{1.2} demonstrará que a integração entre recomendação de
manejo e estimativa de aptidão agrícola, a partir de dados individuais
de solo, representa uma contribuição concreta do SIRAS em relação às
ferramentas atualmente disponíveis. As definições operacionais de
concordância, o conjunto de casos de teste e o procedimento de decisão
sobre as hipóteses são detalhados no \cref{chap:metodologia}.
